\hypertarget{the-future-c26-preview}{%
\chapter{THE FUTURE C26 PREVIEW}\label{the-future-c26-preview}}

\hypertarget{the-future---c26-preview}{%
\chapter{THE FUTURE - C++26 PREVIEW}\label{the-future---c26-preview}}

As of 2026, the C++26 standard is nearing finalization. Here are the
transformative features likely to be included.

\hypertarget{static-reflection-stdmeta}{%
\subsection{13.1 Static Reflection
(std::meta)}\label{static-reflection-stdmeta}}

Reflection allows a program to inspect and modify itself at
compile-time. This eliminates the need for external code generators or
macros for serialization, ORMs, and enum-to-string conversions.

\begin{lstlisting}[language={C++}]
#include <meta>
#include <iostream>
#include <string_view>

struct Person {
    std::string name;
    int age;
    double salary;
};

// Generic serialization using C++26 Reflection
template<typename T>
void serialize(const T& obj) {
    constexpr auto type_info = ^T; // Reflection operator
    
    template for (constexpr auto member : std::meta::members_of(type_info)) {
        std::cout << std::meta::name_of(member) << ": " 
                  << obj.[:member:] << "\n"; // Splicing
    }
}

int main() {
    Person p{"Alice", 30, 95000.0};
    serialize(p); 
    // Output:
    // name: Alice
    // age: 30
    // salary: 95000
}
\end{lstlisting}

\hypertarget{contracts}{%
\subsection{13.2 Contracts}\label{contracts}}

Contracts provide a standardized way to specify preconditions,
postconditions, and assertions, improving safety and optimizer
information.

\begin{lstlisting}[language={C++}]
// pre: Precondition (Caller must ensure)
// post: Postcondition (Function ensures upon return)
// assert: Internal check

int safe_divide(int a, int b) 
    pre { b != 0 }             // Contract: b must not be zero
    post(r) { r * b == a }     // Contract: result * divisor equals dividend
{
    return a / b;
}

// Modes:
// - enforce: Terminate if violated
// - observe: Log/Debug but continue
// - ignore: Optimizer hint (assume true)
\end{lstlisting}

\hypertarget{senders-receivers-stdexecution}{%
\subsection{13.3 Senders \& Receivers
(std::execution)}\label{senders-receivers-stdexecution}}

A unified framework for asynchronous execution, replacing raw threads,
futures, and callbacks with a composable pipeline model.

\begin{lstlisting}[language={C++}]
#include <execution>
#include <iostream>

using namespace std::execution;

int main() {
    scheduler auto sch = thread_pool_scheduler{};

    sender auto work = schedule(sch)
        | then([]{ return 42; })
        | then([](int i){ return i * 2; })
        | then([](int i){ std::cout << "Result: " << i << "\n"; });

    // Launch execution
    std::this_thread::sync_wait(std::move(work));
    
    return 0;
}
\end{lstlisting}

\hypertarget{linear-algebra-stdlinalg}{%
\subsection{13.4 Linear Algebra
(std::linalg)}\label{linear-algebra-stdlinalg}}

Standardized BLAS (Basic Linear Algebra Subprograms) support for
high-performance math.

\begin{lstlisting}[language={C++}]
#include <linalg>
#include <mdspan>
#include <vector>

int main() {
    std::vector<double> A_vec(9), B_vec(3), C_vec(3);
    // ... fill vectors ...

    std::mdspan A(A_vec.data(), 3, 3);
    std::mdspan B(B_vec.data(), 3);
    std::mdspan C(C_vec.data(), 3);

    // Matrix-Vector Multiplication: C = A * B
    std::linalg::matrix_vector_product(A, B, C);
    
    return 0;
}
\end{lstlisting}

\begin{center}\rule{0.5\linewidth}{0.5pt}\end{center}

\begin{center}\rule{0.5\linewidth}{0.5pt}\end{center}

\hypertarget{volume-07-advanced-systems}{%
\chapter{VOLUME 07 ADVANCED SYSTEMS}\label{volume-07-advanced-systems}}
